\documentclass[12pt, titlepage]{article}

\usepackage{fullpage}
\usepackage[round]{natbib}
\usepackage{multirow}
\usepackage{booktabs}
\usepackage{tabularx}
\usepackage{graphicx}
\usepackage{float}
\usepackage{hyperref}
\hypersetup{
    colorlinks,
    citecolor=blue,
    filecolor=black,
    linkcolor=red,
    urlcolor=blue
}

\input{../../Comments.text}
\input{../../Common.text}

\newcounter{acnum}
\newcommand{\actheacnum}{AC\theacnum}
\newcommand{\acref}[1]{AC\ref{#1}}

\newcounter{ucnum}
\newcommand{\uctheucnum}{UC\theucnum}
\newcommand{\uref}[1]{UC\ref{#1}}

\newcounter{mnum}
\newcommand{\mthemnum}{M\themnum}
\newcommand{\mref}[1]{M\ref{#1}}

\begin{document}

\title{Module Guide for \progname{}} 
\author{\authname}
\date{\today}

\maketitle

\pagenumbering{roman}

\section{Revision History}

\begin{tabularx}{\textwidth}{p{3cm}p{2cm}X}
\toprule {\bf Date} & {\bf Version} & {\bf Notes}\\
\midrule
Date 1 & 1.0 & Notes\\
Date 2 & 1.1 & Notes\\
\bottomrule
\end{tabularx}

\newpage

\section{Reference Material}

This section records information for easy reference. 

\subsection{Abbreviations and Acronyms}

\renewcommand{\arraystretch}{1.2}
\begin{tabular}{l l} 
  \toprule		
  \textbf{symbol} & \textbf{description}\\
  \midrule 
  AC & Anticipated Change\\
  DAG & Directed Acyclic Graph \\
  M & Module \\
  MG & Module Guide \\
  OS & Operating System \\
  R & Requirement\\
  SC & Scientific Computing \\
  SRS & Software Requirements Specification\\
  \progname & Explanation of program name\\
  UC & Unlikely Change \\
  \wss{etc.} & \wss{...}\\
  \bottomrule
\end{tabular}\\

\newpage

\tableofcontents

\listoftables

\listoffigures

\newpage

\pagenumbering{arabic}

\section{Introduction}

Decomposing a system into modules is a commonly accepted approach to developing
software.  A module is a work assignment for a programmer or programming
team~\citep{ParnasEtAl1984}.  We advocate a decomposition
based on the principle of information hiding~\citep{Parnas1972a}.  This
principle supports design for change, because the ``secrets'' that each module
hides represent likely future changes.  Design for change is valuable in SC,
where modifications are frequent, especially during initial development as the
solution space is explored.  

Our design follows the rules layed out by \citet{ParnasEtAl1984}, as follows:
\begin{itemize}
\item System details that are likely to change independently should be the
  secrets of separate modules.
\item Each data structure is implemented in only one module.
\item Any other program that requires information stored in a module's data
  structures must obtain it by calling access programs belonging to that module.
\end{itemize}

After completing the first stage of the design, the Software Requirements
Specification (SRS), the Module Guide (MG) is developed~\citep{ParnasEtAl1984}. The MG
specifies the modular structure of the system and is intended to allow both
designers and maintainers to easily identify the parts of the software.  The
potential readers of this document are as follows:

\begin{itemize}
\item New project members: This document can be a guide for a new project member
  to easily understand the overall structure and quickly find the
  relevant modules they are searching for.
\item Maintainers: The hierarchical structure of the module guide improves the
  maintainers' understanding when they need to make changes to the system. It is
  important for a maintainer to update the relevant sections of the document
  after changes have been made.
\item Designers: Once the module guide has been written, it can be used to
  check for consistency, feasibility, and flexibility. Designers can verify the
  system in various ways, such as consistency among modules, feasibility of the
  decomposition, and flexibility of the design.
\end{itemize}

The rest of the document is organized as follows. Section
\ref{SecChange} lists the anticipated and unlikely changes of the software
requirements. Section \ref{SecMH} summarizes the module decomposition that
was constructed according to the likely changes. Section \ref{SecConnection}
specifies the connections between the software requirements and the
modules. Section \ref{SecMD} gives a detailed description of the
modules. Section \ref{SecTM} includes two traceability matrices. One checks
the completeness of the design against the requirements provided in the SRS. The
other shows the relation between anticipated changes and the modules. Section
\ref{SecUse} describes the use relation between modules.

\section{Anticipated and Unlikely Changes} \label{SecChange}

This section lists possible changes to the system. According to the likeliness
of the change, the possible changes are classified into two
categories. Anticipated changes are listed in Section \ref{SecAchange}, and
unlikely changes are listed in Section \ref{SecUchange}.

\subsection{Anticipated Changes} \label{SecAchange}

Anticipated changes are the source of the information that is to be hidden
inside the modules. Ideally, changing one of the anticipated changes will only
require changing the one module that hides the associated decision. The approach
adapted here is called design for
change.

\begin{description}
\item[\refstepcounter{acnum} \actheacnum \label{acHardware}:] The specific
  hardware on which the software is running.
\item[\refstepcounter{acnum} \actheacnum \label{acInput}:] The format of the
  initial input data.
\item ...
\end{description}

\wss{Anticipated changes relate to changes that would be made in requirements,
design or implementation choices.  They are not related to changes that are made
at run-time, like the values of parameters.}

\subsection{Unlikely Changes} \label{SecUchange}

The module design should be as general as possible. However, a general system is
more complex. Sometimes this complexity is not necessary. Fixing some design
decisions at the system architecture stage can simplify the software design. If
these decision should later need to be changed, then many parts of the design
will potentially need to be modified. Hence, it is not intended that these
decisions will be changed.

\begin{description}
\item[\refstepcounter{ucnum} \uctheucnum \label{ucIO}:] Input/Output devices
  (Input: File and/or Keyboard, Output: File, Memory, and/or Screen).
\item ...
\end{description}

\section{Module Hierarchy} \label{SecMH}

This section provides an overview of the module design. Modules are summarized
in a hierarchy decomposed by secrets in Table \ref{TblMH}. The modules listed
below, which are leaves in the hierarchy tree, are the modules that will
actually be implemented.
\\ \\ \indent
\textbf{User Interface Modules}
\begin{description}
\item [\refstepcounter{mnum} \mthemnum \label{mTSI}:] Transcription Service Interface Module
\item \item [\refstepcounter{mnum} \mthemnum \label{mFE}:] Frontend Module
\end{description}
\indent \indent \textbf{Backend Modules}
\begin{description}
\item [\refstepcounter{mnum} \mthemnum \label{mDIO}:] Device I/O Module
\item [\refstepcounter{mnum} \mthemnum \label{mAPI}:] API Gateway Module
\item [\refstepcounter{mnum} \mthemnum \label{mAS}:] Account System Module
\item [\refstepcounter{mnum} \mthemnum \label{mVR}:] Video Recording Module
\item [\refstepcounter{mnum} \mthemnum \label{mRG}:] Report Generation Module
\item [\refstepcounter{mnum} \mthemnum \label{mAR}:] Assessment Review Module
\item [\refstepcounter{mnum} \mthemnum \label{mDB}:] Database Module
\item [\refstepcounter{mnum} \mthemnum \label{mEnc}:] Encryption Module
\item [\refstepcounter{mnum} \mthemnum \label{mML}:] ML Transcription Module
\item [\refstepcounter{mnum} \mthemnum \label{mMERLS}:] MERLS Module
\end{description}

\begin{table}[h!]
\centering
\begin{tabular}{p{0.3\textwidth} p{0.6\textwidth}}
\toprule
\textbf{Level 1} & \textbf{Level 2}\\
\midrule

{Hardware-Hiding} & Device I/O Module \\
\midrule

\multirow{7}{0.3\textwidth}{Behaviour-Hiding} & Frontend Module\\
& Account System Module\\
& Video Recording Module\\
& Transcription Service Interface Module\\
& Report Generation Module\\
& Assessment Review Module\\
& MERLS Module\\
\midrule

\multirow{4}{0.3\textwidth}{Software Decision} & API Gateway Module\\
& Database Module\\
& Encryption Module\\
& ML Transcription Model Module\\
\bottomrule

\end{tabular}
\caption{Module Hierarchy}
\label{TblMH}
\end{table}

\section{Connection Between Requirements and Design} \label{SecConnection}

The design of the system is intended to satisfy the requirements developed in
the SRS. In this stage, the system is decomposed into modules. The connection
between requirements and modules is listed in Table~\ref{TblRT}.

\wss{The intention of this section is to document decisions that are made
  ``between'' the requirements and the design.  To satisfy some requirements,
  design decisions need to be made.  Rather than make these decisions implicit,
  they are explicitly recorded here.  For instance, if a program has security
  requirements, a specific design decision may be made to satisfy those
  requirements with a password.}

\section{Module Decomposition} \label{SecMD}

Modules are decomposed according to the principle of ``information hiding''
proposed by \citet{ParnasEtAl1984}. The \emph{Secrets} field in a module
decomposition is a brief statement of the design decision hidden by the
module. The \emph{Services} field specifies \emph{what} the module will do
without documenting \emph{how} to do it. For each module, a suggestion for the
implementing software is given under the \emph{Implemented By} title. If the
entry is \emph{OS}, this means that the module is provided by the operating
system or by standard programming language libraries.  \emph{\progname{}} means the
module will be implemented by the \progname{} software.

Only the leaf modules in the hierarchy have to be implemented. If a dash
(\emph{--}) is shown, this means that the module is not a leaf and will not have
to be implemented.

\subsection{Hardware Hiding Modules (\mref{mHH})}

\begin{description}
\item[Secrets:]The data structure and algorithm used to implement the virtual
  hardware.
\item[Services:]Serves as a virtual hardware used by the rest of the
  system. This module provides the interface between the hardware and the
  software. So, the system can use it to display outputs or to accept inputs.
\item[Implemented By:] OS
\end{description}

\subsection{Behaviour-Hiding Module}

\begin{description}
\item[Secrets:]The contents of the required behaviours.
\item[Services:]Includes programs that provide externally visible behaviour of
  the system as specified in the software requirements specification (SRS)
  documents. This module serves as a communication layer between the
  hardware-hiding module and the software decision module. The programs in this
  module will need to change if there are changes in the SRS.
\item[Implemented By:] --
\end{description}





\subsubsection{Input Format Module (\mref{mInput})}

\begin{description} %REMOVE WHEN WE SUBMIT
\item[Secrets:]The format and structure of the input data.
\item[Services:]Converts the input data into the data structure used by the
  input parameters module.
\item[Implemented By:] [Your Program Name Here]
\item[Type of Module:] [Record, Library, Abstract Object, or Abstract Data Type]
  [Information to include for leaf modules in the decomposition by secrets tree.]
\end{description}

\subsubsection{Video Recording Module (\mref{mVR})}
\begin{description}
\item[Secrets:]How the recording quality is judged, how a recording moves from being captured to saved, and how failed uploads are retried and verified.
\item[Services:]Checks that the audio and video are good enough before an assessment, records the session, and saves the recording so it can be played back and backed up.
\item[Implemented By:] \progname{}
\item[Type of Module:] Abstract Object
\end{description}

\subsubsection{Transcription Service Interface Module (\mref{mTSI})}
\begin{description}
\item[Secrets:]The format of a transcript, the rules for editing, labelling and verifying transcript segments, and the format of exported annotations.
\item[Services:]Generates a draft transcript for an assessment recording, and lets the clinician correct the text, assign speaker roles, verify segments, mark the transcript complete and export the annotations. 
\item[Type of Module:] Library
\end{description}


\subsection{Software Decision Module}

\begin{description}
\item[Secrets:] The design decision based on mathematical theorems, physical
  facts, or programming considerations. The secrets of this module are
  \emph{not} described in the SRS.
\item[Services:] Includes data structure and algorithms used in the system that
  do not provide direct interaction with the user. 
  % Changes in these modules are more likely to be motivated by a desire to
  % improve performance than by externally imposed changes.
\item[Implemented By:] --
\end{description}

\subsubsection{ML Transcription Model Module (\mref{mML})}
\begin{description}
\item[Secrets:]The speech recognition model used, how audio is split into chunks, and how the model's confidence and the audio level are turned into flags.
\item[Services:]Converts an audio file into timed text segments, each with a confidence score, a detected language and a flag showing whether it needs review.
\item[Implemented By:] \progname{} (using OpenAI's Whisper speech-to-text model)
\item[Type of Module:] Library
\end{description}



\section{Traceability Matrix} \label{SecTM}

This section shows two traceability matrices: between the modules and the
requirements and between the modules and the anticipated changes.

% the table should use mref, the requirements should be named, use something
% like fref
\begin{table}[H]
\centering
\begin{tabular}{p{0.2\textwidth} p{0.6\textwidth}}
\toprule
\textbf{Req.} & \textbf{Modules}\\
\midrule
FR1 & \mref{mTSI}, \mref{mML}\\
FR2 & \mref{mTSI}\\
FR3 & \mref{mTSI}\\
FR6 & \mref{mTSI}\\
FR15 & \mref{mTSI}\\
FR19 & \mref{mTSI}\\
NFR7 & \mref{mVR}, \mref{mTSI}, \mref{mML}\\
NFR8 & \mref{mTSI}\\
NFR9 & \mref{mVR}\\
NFR10 & \mref{mVR}, \mref{mML}\\
NFR13 & \mref{mTSI}\\
NFR14 & \mref{mTSI}\\
NFR17 & \mref{mTSI}, \mref{mML}\\
NFR18 & \mref{mTSI}\\
NFR21 & \mref{mTSI}\\
HA-VR1 & \mref{mVR}\\
HA-VR2 & \mref{mVR}\\
HAMLTM1 & \mref{mML}\\
HAMLTM2 & \mref{mML}\\
\bottomrule
\end{tabular}
\caption{Trace Between Requirements and Modules}
\label{TblRT}
\end{table}

\begin{table}[H]
\centering
\begin{tabular}{p{0.2\textwidth} p{0.6\textwidth}}
\toprule
\textbf{AC} & \textbf{Modules}\\
\midrule
\acref{acHardware} & \mref{mHH}\\
\acref{acInput} & \mref{mInput}\\
\acref{acParams} & \mref{mParams}\\
\acref{acVerify} & \mref{mVerify}\\
\acref{acOutput} & \mref{mOutput}\\
\acref{acVerifyOut} & \mref{mVerifyOut}\\
\acref{acODEs} & \mref{mODEs}\\
\acref{acEnergy} & \mref{mEnergy}\\
\acref{acControl} & \mref{mControl}\\
\acref{acSeqDS} & \mref{mSeqDS}\\
\acref{acSolver} & \mref{mSolver}\\
\acref{acPlot} & \mref{mPlot}\\
\bottomrule
\end{tabular}
\caption{Trace Between Anticipated Changes and Modules}
\label{TblACT}
\end{table}

\section{Use Hierarchy Between Modules} \label{SecUse}

In this section, the uses hierarchy between modules is
provided. \citet{Parnas1978} said of two programs A and B that A {\em uses} B if
correct execution of B may be necessary for A to complete the task described in
its specification. That is, A {\em uses} B if there exist situations in which
the correct functioning of A depends upon the availability of a correct
implementation of B.  Figure \ref{FigUH} illustrates the use relation between
the modules. It can be seen that the graph is a directed acyclic graph
(DAG). Each level of the hierarchy offers a testable and usable subset of the
system, and modules in the higher level of the hierarchy are essentially simpler
because they use modules from the lower levels.

\wss{The uses relation is not a data flow diagram.  In the code there will often
be an import statement in module A when it directly uses module B.  Module B
provides the services that module A needs.  The code for module A needs to be
able to see these services (hence the import statement).  Since the uses
relation is transitive, there is a use relation without an import, but the
arrows in the diagram typically correspond to the presence of import statement.}

\wss{If module A uses module B, the arrow is directed from A to B.}

\begin{figure}[H]
\centering
%\includegraphics[width=0.7\textwidth]{UsesHierarchy.png}
\caption{Use hierarchy among modules}
\label{FigUH}
\end{figure}

%\section*{References}

\section{User Interfaces}

\wss{Design of user interface for software and hardware.  Attach an appendix if
needed. Drawings, Sketches, Figma}

\section{Design of Communication Protocols}

\wss{If there is communication between hardware and software, it may be necessary 
to explain the design of the communication protocol.  This section should be filled 
in with NA if it is not appropriate for a given project.}

\section{Database Design}

\wss{If a database is a significant component of the project, the database
design should be summarized here.  This section should be filled in with NA if
it is not appropriate for a given project.}

\section{Timeline}

\subsection{Module Implementation Timeline (Starting Jan 4)}
This section outlines the implementation timeline for each module. For every module, it describes the module's purpose, identifies the team member responsible for implementing it, and gives the expected start and end dates.

\subsubsection{User Interface (UI) Modules}
\begin{itemize}
    \item \textbf{Jan 7 -- Jan 9: Transcription Service Interface Module (\ref{mTSI})}
    \begin{itemize}
        \item Implement the UI for SAFA in the existing MERLS platform.
        \item Add functionality to interact with the SAFA interface and transcribe.
        \item \textbf{Responsibility:} Nour.
    \end{itemize}

    \item \textbf{Jan 25 -- Jan 27: Frontend Module (\ref{mFE})}
    \begin{itemize}
        \item Implement the clinician frontend.
        \item \textbf{Responsibility:} Catherine.
    \end{itemize}
\end{itemize}

\subsubsection{Backend Modules}
\begin{itemize}
    \item \textbf{Jan 4 -- Jan 21: Device I/O Module (\ref{mDIO})}
    \begin{itemize}
        \item Add functionality to handle hardware input.
        \item \textbf{Responsibility:} Yilin.
    \end{itemize}

    \item \textbf{Jan 4 -- Jan 6: Database Module (\ref{mDB})}
    \begin{itemize}
        \item Implement the database that stores all assessment and patient information.
        \item \textbf{Responsibility:} Javayria.
    \end{itemize}

    \item \textbf{Jan 4 -- Jan 6: ML Transcription Module (\ref{mML})}
    \begin{itemize}
        \item Integrate the existing SAFA tool (SyncScript V3) into MERLS.
        \item Possibly incorporate multiple views for different types of users.
        \item \textbf{Responsibility:} Nour.
    \end{itemize}

    \item \textbf{Jan 7 -- Jan 9: Encryption Module (\ref{mEnc})}
    \begin{itemize}
        \item Implement encryption of patient information.
        \item \textbf{Responsibility:} Javayria.
    \end{itemize}

    \item \textbf{Jan 10 -- Jan 12: Report Generation Module (\ref{mRG})}
    \begin{itemize}
        \item Add functionality to create reports.
        \item \textbf{Responsibility:} Yilin.
    \end{itemize}

    \item \textbf{Jan 13 -- Jan 15: Assessment Review Module (\ref{mAR})}
    \begin{itemize}
        \item Add functionality to flag, save, delete, and mark assessment questions, including creating and deleting notes for questions.
        \item Add functionality to record clinician evaluations and retrieve reviews for patient assessments.
        \item \textbf{Responsibility:} Yilin.
    \end{itemize}

    \item \textbf{Jan 16 -- Jan 18: Account System Module (\ref{mAS})}
    \begin{itemize}
        \item Implement user accounts.
        \item Add functionality for logging in/out and authorizing clinician accounts.
        \item \textbf{Responsibility:} Catherine.
    \end{itemize}

    \item \textbf{Jan 19 -- Jan 21: API Gateway Module (\ref{mAPI})}
    \begin{itemize}
        \item Develop an access point for modules to make their corresponding API calls.
        \item \textbf{Responsibility:} Catherine.
    \end{itemize}

    \item \textbf{Jan 22 -- Jan 24: Video Recording Module (\ref{mVR})}
    \begin{itemize}
        \item Implement a module to handle user recording input (audio and video) for use by the system.
        \item \textbf{Responsibility:} Nour.
    \end{itemize}

    \item \textbf{Jan 28 -- Jan 31: MERLS Module (\ref{mMERLS})}
    \begin{itemize}
        \item Connect the clinician platform to the existing MERLS patient interface.
        \item \textbf{Responsibility:} Javayria.
    \end{itemize}
\end{itemize}

\section{Generative AI Use Disclosure}
This section records all uses of generative artificial intelligence throughout the document.

\subsection{AI Tools Used}

\begin{itemize}
  \item Claude Sonnet 5 (Medium)
  \item ChatGPT-5.6 Luna
\end{itemize}

\subsection{How and Where Used}

The following lists the uses of AI applied throughout the entire document:
\begin{itemize}
  \item Fixing grammar and spelling
  \item Debugging possible \LaTeX{} errors
  \item Performing a final check over the document at the end, comparing each section against the original document comments and the rubric to ensure nothing was missed
  \item Removing template comments
\end{itemize}

\noindent\textbf{Section~\ref{SecMD}}
\begin{itemize}
    \item Check if secrets and services are consistent with the MIS document
\end{itemize}

\noindent\textbf{Section~\ref{SecTM}}
\begin{itemize}
    \item Format both tables for \LaTeX{}
\end{itemize}




\subsection{Verification of AI-Generated Content}

To check the LLM's output, each team member went through their assigned section and presented it to the group, outlining the progress they had made since the previous meeting.
ALL members were able to explain their section with confidence.
Prior to approving any pull request, a team member other than the section's author reviewed it to make sure it was cohesive, accurate, and in consistent with the rest of the document.
Since the document did not include any factual claims, the LLM's output did not need further fact-checking.

\newpage{}

\bibliographystyle {plainnat}
\bibliography{../../../refs/References}

\end{document}